\begin{derivation}[从\eqnrefs{eq:qcd-gluon-propagator-dp10}到\eqnrefs{eq:qcd-three-gluon-vertex-dp10}]
Feynman 规范 $\xi=1$ 下，把纯规范二次项与规范固定项分部积分后写成
\begin{equation*}
 \mathcal L_2
 =\frac{\hbar c}{2}\,
 \mathscr C_\mu^a\,\eta^{\mu\nu}\partial^2\mathscr C_\nu^a .
\end{equation*}
因此作用在规范场上的二次微分算符是
\begin{equation*}
 (\mathsf K\mathscr C)_\mu^a
 =\delta^{ab}\eta_{\mu\nu}\partial^2\mathscr C^{b\nu}.
\end{equation*}
采用
\begin{equation*}
 \mathscr C_\mu^a(x)
 =\int\frac{\dd^4p}{(2\pi\hbar)^4}
 \widetilde{\mathscr C}_\mu^a(p)
 \ee^{-\ii p\cdot x/\hbar},
\end{equation*}
有 $\partial^2\mapsto-p^2/\hbar^2$，故
\begin{equation*}
 \mathsf K_{\mu\nu}^{ab}(p)
 =-\frac{p^2}{\hbar^2}\delta^{ab}\eta_{\mu\nu}.
\end{equation*}
Feynman 边界条件把 $p^2$ 替换为 $p^2+\ii0^+$。定义逆核 $\mathsf K_F^{-1}$ 使
\begin{equation*}
 \mathsf K_{\mu\rho}^{ac}(p)
 (\mathsf K_F^{-1})^{cb\rho}{}_{\nu}(p)
 =\delta^{ab}\delta_\mu{}^\nu,
\end{equation*}
于是逐指标求逆得到
\begin{equation*}
 (\mathsf K_F^{-1})_{\mu\nu}^{ab}(p)
 =-\hbar^2\frac{\delta^{ab}\eta_{\mu\nu}}{p^2+\ii0^+}.
\end{equation*}
在统一的约化振幅约定中，玻色 Feynman 线等于 $(\ii/\hbar^2)\mathsf K_F^{-1}$，从而
\begin{equation*}
 D_{\mu\nu}^{ab}(p)
 =\frac{\ii}{\hbar^2}(\mathsf K_F^{-1})_{\mu\nu}^{ab}(p)
 =\frac{-\ii\delta^{ab}\eta_{\mu\nu}}{p^2+\ii0^+},
\end{equation*}
即正文\eqnrefs{eq:qcd-gluon-propagator-dp10}。

三胶子规则来自 Yang--Mills 作用量的三次部分
\begin{equation*}
 \mathcal L_{3g}
 =\hbar c g_s f^{abc}
 (\partial_\mu\mathscr C_\nu^a)
 \mathscr C^{b\mu}\mathscr C^{c\nu}.
\end{equation*}
令三条流入场分别为 $(a,\alpha,p)$、$(b,\beta,q)$、$(c,\gamma,r)$，并约定
\begin{equation*}
 p+q+r=0.
\end{equation*}
三次场都是同一个玻色规范场，因此对作用量取三次泛函导数时，导数场可以对应三条外腿中的任意一条。把所有排列写到统一的颜色次序 $f^{abc}$ 后，三个 Lorentz 系数分别为
\begin{align*}
 \mathscr V_1^{\alpha\beta\gamma}
 &=\eta^{\alpha\beta}(p-q)^\gamma,\\
 \mathscr V_2^{\alpha\beta\gamma}
 &=\eta^{\beta\gamma}(q-r)^\alpha,\\
 \mathscr V_3^{\alpha\beta\gamma}
 &=\eta^{\gamma\alpha}(r-p)^\beta.
\end{align*}
这里每一个动量差都来自两种外腿排列的差：交换相邻两个胶子场时，$f^{abc}$ 变号，而 Fourier 导数把相应流入动量带入同一个 Lorentz 缩并。于是三阶作用量系数给出完整约化顶角
\begin{equation*}
 \boxed{
 V_{3g}^{abc,\alpha\beta\gamma}(p,q,r)
 =-g_sf^{abc}
 \left[
 \eta^{\alpha\beta}(p-q)^\gamma
 +\eta^{\beta\gamma}(q-r)^\alpha
 +\eta^{\gamma\alpha}(r-p)^\beta
 \right],}
\end{equation*}
即正文\eqnrefs{eq:qcd-three-gluon-vertex-dp10}。交换任意两条完整外腿，例如 $(a,\alpha,p)\leftrightarrow(b,\beta,q)$，$f^{abc}$ 与方括号同时变号，因此完整顶角保持玻色交换对称。
\end{derivation}
